\subsection{INSERT (I1--I6)}

Grupa scenariuszy INSERT obejmuje sześć wariantów wstawiania danych -- od pojedynczego
rekordu z powiązaniami, przez masowe ładowanie paczek rekordów, aż po wstawianie złożonych struktur
hierarchicznych. Scenariusze różnią się
wolumenem wstawianych danych oraz stopniem powiązań między rekordami, co pozwala
porównać efektywność mechanizmów masowego ładowania w każdym z silników (COPY,
multi-row INSERT, \texttt{insertMany}, \texttt{UNWIND + CALL IN TRANSACTIONS}) oraz
koszt utrzymania spójności relacyjnej i referencyjnej przy zapisie.

\subsubsection{I1: Rejestracja użytkownika}

Rejestracja nowego użytkownika platformy -- w ramach jednej operacji tworzony
jest rekord użytkownika, jego domyślny profil oraz aktywna subskrypcja w planie
Standard. Scenariusz testuje koszt wstawienia powiązanych rekordów w trzech
różnych tabelach/kolekcjach/węzłach grafu oraz utrzymania ich spójności
(relacja 1-1 users→subscription, 1-N users→profiles).

\textbf{PostgreSQL}
\begin{lstlisting}[language=SQL]
INSERT INTO users (email, password_hash, first_name, last_name,
    date_of_birth, country_code, status, created_at, updated_at)
VALUES (%s, '$2b$12$benchmark', 'Bench', 'User',
    '1990-01-01', 'PL', 'active', NOW(), NOW())
RETURNING user_id;

INSERT INTO profiles (user_id, name, maturity_rating, language, created_at)
VALUES (%s, 'Main', 'ALL', 'pl', NOW());

INSERT INTO subscriptions (user_id, plan_name, price_monthly, max_streams,
    max_resolution, status, start_date, auto_renew, created_at)
VALUES (%s, 'Standard', 43.99, 2, '1080p', 'active', CURRENT_DATE, TRUE, NOW());
\end{lstlisting}

\textbf{MySQL}
\begin{lstlisting}[language=SQL]
INSERT INTO users (email, password_hash, first_name, last_name,
    date_of_birth, country_code, status, created_at, updated_at)
VALUES (%s, '$2b$12$benchmark', 'Bench', 'User',
    '1990-01-01', 'PL', 'active', NOW(), NOW());

INSERT INTO profiles (user_id, name, maturity_rating, language, created_at)
VALUES (%s, 'Main', 'ALL', 'pl', NOW());

INSERT INTO subscriptions (user_id, plan_name, price_monthly, max_streams,
    max_resolution, status, start_date, auto_renew, created_at)
VALUES (%s, 'Standard', 43.99, 2, '1080p', 'active', CURDATE(), TRUE, NOW());
\end{lstlisting}

\textbf{MongoDB}
\begin{lstlisting}[language=Python]
db.users.insertOne({
    _id: <uid>,
    email: <email>, password_hash: "$2b$12$benchmark",
    first_name: "Bench", last_name: "User",
    date_of_birth: "1990-01-01", country_code: "PL",
    status: "active",
    created_at: "2025-01-01 00:00:00",
    updated_at: "2025-01-01 00:00:00",
    profiles: [{
        profile_id: <uid * 10>, name: "Main", is_kids: false,
        maturity_rating: "ALL", language: "pl",
        created_at: "2025-01-01 00:00:00"
    }],
    subscription: {
        subscription_id: <uid * 10>, plan_name: "Standard",
        price_monthly: 43.99, max_streams: 2, max_resolution: "1080p",
        status: "active", start_date: "2025-01-01", auto_renew: true
    }
});
\end{lstlisting}

\textbf{Neo4j}
\begin{lstlisting}[language=SQL]
CREATE (u:User {user_id: $uid, email: $email,
    first_name: 'Bench', last_name: 'User',
    country_code: 'PL', status: 'active',
    created_at: '2025-01-01 00:00:00'})
CREATE (p:Profile {profile_id: $pid, name: 'Main',
    is_kids: false, maturity_rating: 'ALL', language: 'pl'})
CREATE (sub:Subscription {subscription_id: $sid,
    plan_name: 'Standard', price_monthly: 43.99,
    max_streams: 2, max_resolution: '1080p',
    status: 'active', start_date: '2025-01-01', auto_renew: true})
CREATE (u)-[:HAS_PROFILE]->(p)
CREATE (u)-[:HAS_SUBSCRIPTION]->(sub);
\end{lstlisting}


\subsubsection{I2: Bulk import ocen}

Masowe wczytanie dużej paczki ocen (10k / 100k / 500k w zależności od wolumenu).
Mierzy najbardziej efektywną ścieżkę ładowania danych w każdej z baz: COPY
w PostgreSQL, wielowierszowy INSERT IGNORE w MySQL, insert\_many w MongoDB
oraz UNWIND + CALL IN TRANSACTIONS w Neo4j. Testuje przepustowość masowego
zapisu z eliminacją powtórzeń po unikalnej parze (profile\_id, content\_id).

\textbf{PostgreSQL}
\begin{lstlisting}[language=SQL]
CREATE TEMP TABLE _stage_ratings (
    profile_id BIGINT, content_id BIGINT, score INT,
    review_text TEXT, created_at TIMESTAMP, updated_at TIMESTAMP
) ON COMMIT DROP;

COPY _stage_ratings (profile_id, content_id, score,
    review_text, created_at, updated_at) FROM STDIN;

INSERT INTO ratings (profile_id, content_id, score,
    review_text, created_at, updated_at)
SELECT profile_id, content_id, score,
    review_text, created_at, updated_at FROM _stage_ratings
ON CONFLICT (profile_id, content_id) DO NOTHING;
\end{lstlisting}

\textbf{MySQL}
\begin{lstlisting}[language=SQL]
-- batch po 5000 wierszy
INSERT IGNORE INTO ratings (profile_id, content_id, score,
    review_text, created_at, updated_at)
VALUES (%s, %s, %s, %s, %s, %s), ... ;
\end{lstlisting}

\textbf{MongoDB}
\begin{lstlisting}[language=Python]
-- batch po 5000 dokumentow
db.ratings.insertMany([
    { _id: <id>, profile_id: <pid>, content_id: <cid>, score: <1..10>,
      review_text: "", created_at: "2037-12-31 00:00:00",
      updated_at: "2037-12-31 00:00:00" }, ...
], { ordered: false });
\end{lstlisting}

\textbf{Neo4j}
\begin{lstlisting}[language=SQL]
-- batch po 25 000, z wewnetrznym TRANSACTIONS OF 2000 ROWS
UNWIND $rows AS r
CALL (r) {
    MATCH (p:Profile {profile_id: r.profile_id})
    MATCH (c:Content {content_id: r.content_id})
    CREATE (p)-[:RATED {score: r.score, rated_at: r.rated_at}]->(c)
} IN TRANSACTIONS OF 2000 ROWS;
\end{lstlisting}


\subsubsection{I3: Dodanie serialu z pełnym drzewem}

Dodanie nowego serialu wraz z pełnym hierarchicznym drzewem zależności:
treść (content) → 3 sezony → po 10 odcinków w każdym sezonie, plus 10 relacji
z osobami (reżyser, scenarzysta, aktorzy). Pokazuje koszt tworzenia silnie
powiązanych, zagnieżdżonych struktur: normalizacja w SQL (4 tabele + tabela
łącząca), zagnieżdżony dokument w Mongo oraz wieloetapowe CREATE + MATCH w Cypher.

\textbf{PostgreSQL}
\begin{lstlisting}[language=SQL]
INSERT INTO content (title, description, type, release_date,
    maturity_rating, genres, avg_rating, popularity_score,
    country_of_origin, original_language, is_active, metadata, created_at)
VALUES ('Benchmark Series', 'Test', 'series', '2025-01-01',
    '16+', 'Drama,Thriller', 0, 0, 'PL', 'pl', TRUE, '{...}'::jsonb, NOW())
RETURNING content_id;

-- dla sezonow 1..3:
INSERT INTO seasons (content_id, season_number, title, release_date)
VALUES (%s, %s, %s, '2025-01-01') RETURNING season_id;

-- dla odcinkow 1..10:
INSERT INTO episodes (season_id, episode_number, title,
    description, duration_minutes, release_date, video_url)
VALUES (%s, %s, %s, 'Test episode', 45, '2025-01-01',
    'https://cdn.example.com/test.mp4');

-- dla 10 osob:
INSERT INTO content_people (content_id, person_id, role,
    character_name, billing_order)
VALUES (%s, %s, %s, %s, %s)
ON CONFLICT DO NOTHING;
\end{lstlisting}

\textbf{MySQL}
\begin{lstlisting}[language=SQL]
INSERT INTO content (...) VALUES ('Benchmark Series', ..., '{...}', NOW());
-- analogicznie jak w PostgreSQL, bez RETURNING (uzywa lastrowid)
INSERT INTO seasons (...) VALUES (...);
INSERT INTO episodes (...) VALUES (...);
INSERT IGNORE INTO content_people (content_id, person_id, role,
    character_name, billing_order) VALUES (%s, %s, %s, %s, %s);
\end{lstlisting}

\textbf{MongoDB}
\begin{lstlisting}[language=Python]
db.content.insertOne({
    _id: <cid>, title: "Benchmark Series", type: "series",
    release_date: "2025-01-01", maturity_rating: "16+",
    genres: ["Drama", "Thriller"], avg_rating: 0, is_active: true,
    cast: [ { person_id, first_name, last_name, role: "actor",
              character_name, billing_order }, ... ],
    seasons: [
        { season_id, season_number, title, release_date,
          episodes: [ { episode_id, episode_number, title,
                        duration_minutes: 45, release_date }, ... ]
        }, ...
    ]
});
\end{lstlisting}

\textbf{Neo4j}
\begin{lstlisting}[language=SQL]
CREATE (c:Content {content_id: $cid, title: 'Benchmark Series',
    type: 'series', maturity_rating: '16+',
    avg_rating: 0, is_active: true, created_at: '2025-01-01 00:00:00'})
WITH c
UNWIND range(1, 3) AS sn
CREATE (s:Season {season_id: $cid * 100 + sn, season_number: sn})
CREATE (c)-[:HAS_SEASON]->(s)
WITH s, sn
UNWIND range(1, 10) AS en
CREATE (e:Episode {episode_id: $cid * 1000 + sn * 100 + en,
    episode_number: en, duration_minutes: 45})
CREATE (s)-[:HAS_EPISODE]->(e);

UNWIND $pids AS pid
MATCH (p:Person {person_id: pid})
MATCH (c:Content {content_id: $cid})
CREATE (p)-[:ACTED_IN {character_name: 'Test',
    billing_order: pid % 10}]->(c);
\end{lstlisting}


\subsubsection{I4: Batch insert płatności}

Hurtowe wstawienie paczki płatności (1k / 10k / 50k) bez sprawdzania
powtórzeń -- czysty pomiar zimnego wstawiania rekordów z prostą relacją
do istniejących subskrypcji. Każda baza używa swojego najszybszego mechanizmu batchowego
(COPY, multi-row INSERT, insert\_many, UNWIND+CREATE).

\textbf{PostgreSQL}
\begin{lstlisting}[language=SQL]
COPY payments (subscription_id, amount, currency,
    payment_method, transaction_id, status, paid_at, created_at)
FROM STDIN;
\end{lstlisting}

\textbf{MySQL}
\begin{lstlisting}[language=SQL]
-- batch po 5000
INSERT INTO payments (subscription_id, amount, currency,
    payment_method, transaction_id, status, paid_at, created_at)
VALUES (%s, %s, %s, %s, %s, %s, %s, %s), ... ;
\end{lstlisting}

\textbf{MongoDB}
\begin{lstlisting}[language=Python]
-- batch po 5000
db.payments.insertMany([
    { _id: <id>, subscription_id, amount, currency: "PLN",
      payment_method, transaction_id, status: "completed",
      paid_at: "2025-06-15 12:00:00",
      created_at: "2025-06-15 12:00:00" }, ...
], { ordered: false });
\end{lstlisting}

\textbf{Neo4j}
\begin{lstlisting}[language=SQL]
-- batch po 5000
UNWIND $rows AS r
MATCH (sub:Subscription {subscription_id: r.subscription_id})
CREATE (pay:Payment {payment_id: r.payment_id,
    amount: r.amount, currency: r.currency,
    payment_method: r.payment_method,
    status: r.status, paid_at: r.paid_at})
CREATE (sub)-[:HAS_PAYMENT]->(pay);
\end{lstlisting}


\subsubsection{I5: Oceny z przeliczeniem średniej}

Kombinowana operacja: wstawienie paczki ocen (500 / 5k / 25k) dla losowo
wybranej treści, po czym przeliczenie i zapisanie zaktualizowanej średniej
avg\_rating tej treści. Mierzy łączny koszt INSERT + agregacja AVG + UPDATE.

\textbf{PostgreSQL}
\begin{lstlisting}[language=SQL]
INSERT INTO ratings (profile_id, content_id, score,
    review_text, created_at, updated_at)
VALUES (%s, %s, %s, %s, %s, %s)
ON CONFLICT (profile_id, content_id) DO NOTHING;

UPDATE content SET avg_rating = COALESCE(
    (SELECT AVG(score) FROM ratings WHERE content_id = %s), 0)
WHERE content_id = %s;
\end{lstlisting}

\textbf{MySQL}
\begin{lstlisting}[language=SQL]
-- batch po 5000
INSERT IGNORE INTO ratings (profile_id, content_id, score,
    review_text, created_at, updated_at)
VALUES (%s, %s, %s, %s, %s, %s), ... ;

UPDATE content SET avg_rating = COALESCE(
    (SELECT AVG(score) FROM ratings WHERE content_id = %s), 0)
WHERE content_id = %s;
\end{lstlisting}

\textbf{MongoDB}
\begin{lstlisting}[language=Python]
db.ratings.insertMany([ ... ], { ordered: false });

db.ratings.aggregate([
    { $match: { content_id: <cid> } },
    { $group: { _id: null, avg: { $avg: "$score" } } }
]);

db.content.updateOne(
    { _id: <cid> },
    { $set: { avg_rating: <avg> } }
);
\end{lstlisting}

\textbf{Neo4j}
\begin{lstlisting}[language=SQL]
-- batch po 5000
UNWIND $rows AS r
MATCH (p:Profile {profile_id: r.profile_id})
MATCH (c:Content {content_id: r.content_id})
CREATE (p)-[:RATED {score: r.score,
    created_at: '2025-06-15 12:00:00'}]->(c);

MATCH (:Profile)-[r:RATED]->(c:Content {content_id: $cid})
WITH c, avg(r.score) AS avgScore
SET c.avg_rating = avgScore;
\end{lstlisting}


\subsubsection{I6: Import osób z powiązaniami}

Import paczki osób (200 / 1k / 5k) wraz z powiązaniem każdej z nich relacją
"actor" z losowo wybraną treścią. Testuje dwuetapowe ładowanie: najpierw same
encje, następnie relacje -- istotne dla grafu (Neo4j wymaga MATCH po istniejących węzłach przed
utworzeniem krawędzi) oraz dla modelu zagnieżdżonego w Mongo (dołączenie
do tablicy cast wewnątrz dokumentu treści).

\textbf{PostgreSQL}
\begin{lstlisting}[language=SQL]
COPY people (person_id, first_name, last_name, birth_date,
    bio, photo_url, nationality) FROM STDIN;

INSERT INTO content_people (content_id, person_id, role,
    character_name, billing_order)
VALUES (%s, %s, %s, %s, %s)
ON CONFLICT DO NOTHING;
\end{lstlisting}

\textbf{MySQL}
\begin{lstlisting}[language=SQL]
-- batch po 5000
INSERT INTO people (person_id, first_name, last_name,
    birth_date, bio, photo_url, nationality)
VALUES (%s, %s, %s, %s, %s, %s, %s), ... ;

-- batch po 5000
INSERT IGNORE INTO content_people (content_id, person_id, role,
    character_name, billing_order)
VALUES (%s, %s, %s, %s, %s), ... ;
\end{lstlisting}

\textbf{MongoDB}
\begin{lstlisting}[language=Python]
db.content.bulkWrite([
    { updateOne: {
        filter: { _id: <cid> },
        update: { $push: { cast: {
            person_id: <pid>, first_name, last_name,
            role: "actor", character_name, billing_order
        } } }
    }},
    ...
], { ordered: false });
\end{lstlisting}

\textbf{Neo4j}
\begin{lstlisting}[language=SQL]
-- batch po 5000: tworzenie osob
UNWIND $rows AS r
CREATE (:Person {person_id: r.person_id,
    first_name: r.first_name, last_name: r.last_name,
    birth_date: r.birth_date, nationality: r.nationality});

-- batch po 5000: relacje
UNWIND $rows AS r
MATCH (p:Person {person_id: r.person_id})
MATCH (c:Content {content_id: r.content_id})
CREATE (p)-[:ACTED_IN {character_name: r.character_name,
    billing_order: r.billing_order}]->(c);
\end{lstlisting}


% -------------------------------------------------------
\subsection{SELECT (S1--S6)}
% -------------------------------------------------------

Grupa scenariuszy SELECT obejmuje typowe zapytania odczytu występujące w platformie
VOD -- od szybkich zapytań listujących z sortowaniem, przez złożone agregacje 
z filtrem zakresu dat oraz  wyszukiwanie pełnotekstowe, aż po zapytania grafowe
silnie obciążające JOIN-y. Zestaw jest dobrany tak, aby pokazać
różnice między bazami w trzech wymiarach: szybkości prostych odczytów po kluczu
głównym, kosztu agregacji i JOIN-ów oraz efektywności indeksów wtórnych i pełnotekstowych.

\subsubsection{S1: Strona główna}

Typowe zapytanie strony głównej serwisu VOD -- zwraca 20 najbardziej popularnych,
aktywnych treści z gatunku Action, posortowanych malejąco po popularity\_score.
Testuje filtrowanie po zagnieżdżonym polu (genres) + sort + limit oraz odczyt
pola z dokumentu JSON (metadata.studio).

\textbf{PostgreSQL}
\begin{lstlisting}[language=SQL]
SELECT content_id, title, genres, popularity_score, avg_rating,
       metadata->>'studio' AS studio
FROM content
WHERE is_active = TRUE AND genres LIKE '%Action%'
ORDER BY popularity_score DESC
LIMIT 20;
\end{lstlisting}

\textbf{MySQL}
\begin{lstlisting}[language=SQL]
SELECT content_id, title, genres, popularity_score, avg_rating,
       metadata->>'$.studio' AS studio
FROM content
WHERE is_active = TRUE AND genres LIKE '%Action%'
ORDER BY popularity_score DESC
LIMIT 20;
\end{lstlisting}

\textbf{MongoDB}
\begin{lstlisting}[language=Python]
db.content.find(
    { is_active: true, genres: "Action" },
    { title: 1, genres: 1, popularity_score: 1,
      avg_rating: 1, "metadata.studio": 1 }
).sort({ popularity_score: -1 }).limit(20);
\end{lstlisting}

\textbf{Neo4j}
\begin{lstlisting}[language=SQL]
MATCH (c:Content)-[:HAS_GENRE]->(g:Genre {name: 'Action'})
WHERE c.is_active = true
RETURN c.content_id, c.title, c.popularity_score, c.avg_rating
ORDER BY c.popularity_score DESC
LIMIT 20;
\end{lstlisting}


\subsubsection{S2: Filtrowanie kolaboratywne}

Rekomendacje w oparciu o filtrowanie kolaboratywne -- znajdź profile oglądające
te same treści co wybrany profil, a następnie zwróć 10 treści najczęściej
oglądanych przez tę podobną grupę, których użytkownik jeszcze nie oglądał.
Scenariusz silnie obciąża JOIN-y po watch\_history i jest naturalnym testem
przewagi grafu (Neo4j).

\textbf{PostgreSQL / MySQL}
\begin{lstlisting}[language=SQL]
WITH similar_profiles AS (
    SELECT DISTINCT wh2.profile_id
    FROM watch_history wh1
    JOIN watch_history wh2
        ON wh1.content_id = wh2.content_id
        AND wh1.profile_id <> wh2.profile_id
    WHERE wh1.profile_id = %s
    LIMIT 50
)
SELECT c.content_id, c.title, COUNT(*) AS score
FROM watch_history wh
JOIN content c ON wh.content_id = c.content_id
WHERE wh.profile_id IN (SELECT profile_id FROM similar_profiles)
  AND wh.content_id NOT IN (
      SELECT content_id FROM watch_history WHERE profile_id = %s
  )
GROUP BY c.content_id, c.title
ORDER BY score DESC
LIMIT 10;
\end{lstlisting}

\textbf{MongoDB}
\begin{lstlisting}[language=Python]
watched = db.watch_history.distinct("content_id", { profile_id: <pid> });

similar = db.watch_history.distinct("profile_id",
    { content_id: { $in: watched[:50] },
      profile_id: { $ne: <pid> } })[:50];

db.watch_history.aggregate([
    { $match: { profile_id: { $in: similar },
                content_id: { $nin: watched } } },
    { $group: { _id: "$content_id", score: { $sum: 1 } } },
    { $sort: { score: -1 } },
    { $limit: 10 },
    { $lookup: { from: "content", localField: "_id",
                 foreignField: "_id", as: "content" } }
]);
\end{lstlisting}

\textbf{Neo4j}
\begin{lstlisting}[language=SQL]
MATCH (p:Profile {profile_id: $pid})-[:WATCHED]->(c:Content)
      <-[:WATCHED]-(similar:Profile)
WHERE similar <> p
WITH p, similar, COUNT(c) AS shared
ORDER BY shared DESC LIMIT 50
MATCH (similar)-[:WATCHED]->(rec:Content)
WHERE NOT EXISTS { (p)-[:WATCHED]->(rec) }
RETURN rec.content_id, rec.title, COUNT(*) AS score
ORDER BY score DESC LIMIT 10;
\end{lstlisting}


\subsubsection{S3: TOP 100 według liczby ocen w okresie}

Ranking 100 treści o największej liczbie ocen wystawionych po określonej dacie
odcięcia, wraz ze średnią oceną. Mierzy agregację GROUP BY z filtrem zakresu
dat oraz JOIN z treścią w celu pobrania tytułu. Klasyczne zapytanie analityczne
typu "co jest ostatnio na topie".

\textbf{PostgreSQL / MySQL}
\begin{lstlisting}[language=SQL]
SELECT c.content_id, c.title, COUNT(*) AS rating_count,
       AVG(r.score) AS avg_score
FROM ratings r
JOIN content c ON r.content_id = c.content_id
WHERE r.created_at >= '2025-06-01'
GROUP BY c.content_id, c.title
ORDER BY rating_count DESC
LIMIT 100;
\end{lstlisting}

\textbf{MongoDB}
\begin{lstlisting}[language=Python]
db.ratings.aggregate([
    { $match: { created_at: { $gte: "2025-06-01" } } },
    { $group: { _id: "$content_id",
                rating_count: { $sum: 1 },
                avg_score: { $avg: "$score" } } },
    { $sort: { rating_count: -1 } },
    { $limit: 100 },
    { $lookup: { from: "content", localField: "_id",
                 foreignField: "_id", as: "content" } }
]);
\end{lstlisting}

\textbf{Neo4j}
\begin{lstlisting}[language=SQL]
MATCH ()-[r:RATED]->(c:Content)
WHERE r.rated_at >= '2025-06-01'
RETURN c.content_id, c.title,
       COUNT(r) AS rating_count, avg(r.score) AS avg_score
ORDER BY rating_count DESC
LIMIT 100;
\end{lstlisting}


\subsubsection{S4: Wyszukiwanie pełnotekstowe po tytule}

Wyszukiwanie treści po fragmencie tytułu. Wariant bez indeksów wymusza pełny
skan, wariant z indeksami korzysta z wyspecjalizowanych struktur pełnotekstowych
(GIN+pg\_trgm, FULLTEXT, \$text, db.index.fulltext) -- pozwala porównać zysk
z indeksu tekstowego między silnikami.

\textbf{PostgreSQL}
\begin{lstlisting}[language=SQL]
SELECT content_id, title, popularity_score,
       metadata->>'tags' AS tags
FROM content
WHERE title ILIKE '%Interface%'
ORDER BY popularity_score DESC
LIMIT 20;
\end{lstlisting}

\textbf{MySQL}
\begin{lstlisting}[language=SQL]
SELECT content_id, title, popularity_score,
       metadata->>'$.tags' AS tags
FROM content
WHERE title LIKE '%Interface%'
ORDER BY popularity_score DESC
LIMIT 20;
\end{lstlisting}

\textbf{MongoDB}
\begin{lstlisting}[language=Python]
-- z indeksem full-text ($text):
db.content.find(
    { $text: { $search: "Interface" } },
    { score: { $meta: "textScore" },
      title: 1, popularity_score: 1, "metadata.tags": 1 }
).sort({ score: { $meta: "textScore" } }).limit(20);

-- bez indeksu (regex):
db.content.find(
    { title: { $regex: /Interface/i } },
    { title: 1, popularity_score: 1, "metadata.tags": 1 }
).sort({ popularity_score: -1 }).limit(20);
\end{lstlisting}

\textbf{Neo4j}
\begin{lstlisting}[language=SQL]
-- z indeksem full-text:
CALL db.index.fulltext.queryNodes('content_title_ft', 'Interface')
YIELD node AS c, score
RETURN c.content_id, c.title, c.popularity_score
ORDER BY c.popularity_score DESC
LIMIT 20;

-- bez indeksu:
MATCH (c:Content)
WHERE c.title CONTAINS 'Interface'
RETURN c.content_id, c.title, c.popularity_score
ORDER BY c.popularity_score DESC
LIMIT 20;
\end{lstlisting}


\subsubsection{S5: Historia oglądania profilu}

Ostatnie 50 pozycji historii oglądania danego profilu, wzbogacone o tytuł i typ
treści. Klasyczne zapytanie ekranu "Kontynuuj oglądanie" -- sort malejąco
po started\_at.

\textbf{PostgreSQL / MySQL}
\begin{lstlisting}[language=SQL]
SELECT wh.watch_id, c.title, c.type, wh.started_at,
       wh.progress_percent, wh.completed
FROM watch_history wh
JOIN content c ON wh.content_id = c.content_id
WHERE wh.profile_id = %s
ORDER BY wh.started_at DESC
LIMIT 50;
\end{lstlisting}

\textbf{MongoDB}
\begin{lstlisting}[language=Python]
db.watch_history.aggregate([
    { $match: { profile_id: <pid> } },
    { $sort: { started_at: -1 } },
    { $limit: 50 },
    { $lookup: { from: "content", localField: "content_id",
                 foreignField: "_id", as: "content" } },
    { $unwind: "$content" },
    { $project: { "content.title": 1, "content.type": 1,
                  started_at: 1, progress_percent: 1, completed: 1 } }
]);
\end{lstlisting}

\textbf{Neo4j}
\begin{lstlisting}[language=SQL]
MATCH (p:Profile {profile_id: $pid})-[w:WATCHED]->(c:Content)
RETURN c.title, c.type, w.started_at,
       w.progress_percent, w.completed
ORDER BY w.started_at DESC
LIMIT 50;
\end{lstlisting}


\subsubsection{S6: Filmografia osoby}

Pełna filmografia wybranej osoby -- wszystkie treści powiązane z nią dowolną
rolą (actor / director / writer), posortowane od najnowszej. W grafie
realizowane przez wzorzec zmiennej etykiety relacji, w Mongo przez indeks
na "cast.person\_id".

\textbf{PostgreSQL / MySQL}
\begin{lstlisting}[language=SQL]
SELECT c.content_id, c.title, c.type, c.release_date,
       cp.role, cp.character_name
FROM content_people cp
JOIN content c ON cp.content_id = c.content_id
WHERE cp.person_id = %s
ORDER BY c.release_date DESC;
\end{lstlisting}

\textbf{MongoDB}
\begin{lstlisting}[language=Python]
db.content.find(
    { "cast.person_id": <person_id> },
    { title: 1, type: 1, release_date: 1, "cast.$": 1 }
).sort({ release_date: -1 });
\end{lstlisting}

\textbf{Neo4j}
\begin{lstlisting}[language=SQL]
MATCH (p:Person {person_id: $pid})-[r]->(c:Content)
WHERE type(r) IN ['ACTED_IN', 'DIRECTED', 'WROTE']
RETURN c.content_id, c.title, c.type, c.release_date,
       type(r) AS role, r.character_name
ORDER BY c.release_date DESC;
\end{lstlisting}


% -------------------------------------------------------
\subsection{UPDATE (U1--U6)}
% -------------------------------------------------------

Grupa scenariuszy UPDATE obejmuje sześć wariantów modyfikacji istniejących danych --
od najprostszej aktualizacji pojedynczego rekordu po kluczu głównym, przez 
przeliczenie pola z agregatu, aż po masowe aktualizacje setek tysięcy wierszy po filtrze
niebędącym kluczem głównym. Zestaw pozwala zweryfikować drugą część hipotezy, czyli
hipotetyczny spadek wydajności UPDATE przy obecności indeksów, a jednocześnie
porównać modele danych w kontekście aktualizacji zagnieżdżonych pól (MongoDB)
oraz właściwości krawędzi grafu (Neo4j).

\subsubsection{U1: Aktualizacja postępu oglądania}

Aktualizacja postępu oglądania pojedynczego wpisu historii (progress\_percent,
completed). Typowa aktualizacja jednego rekordu po kluczu głównym -- mierzy czysty koszt
dotarcia do rekordu i zapisu pary pól bez żadnych dodatkowych zależności.

\textbf{PostgreSQL / MySQL}
\begin{lstlisting}[language=SQL]
UPDATE watch_history
SET progress_percent = 75.50, completed = FALSE
WHERE watch_id = %s;
\end{lstlisting}

\textbf{MongoDB}
\begin{lstlisting}[language=Python]
db.watch_history.updateOne(
    { _id: <watch_id> },
    { $set: { progress_percent: 75.50, completed: false } }
);
\end{lstlisting}

\textbf{Neo4j}
\begin{lstlisting}[language=SQL]
MATCH (p:Profile {profile_id: $pid})-[w:WATCHED]->(c:Content)
WITH w ORDER BY w.started_at DESC LIMIT 1
SET w.progress_percent = 75.50, w.completed = false;
\end{lstlisting}


\subsubsection{U2: Przeliczenie średniej oceny}

Przeliczenie i zapis średniej oceny (avg\_rating) dla jednej treści na podstawie
wszystkich rekordów tabeli ratings przypisanych do tej treści. Mierzy koszt
agregacji AVG + pojedynczego UPDATE dla różnych modeli danych.

\textbf{PostgreSQL}
\begin{lstlisting}[language=SQL]
UPDATE content SET avg_rating = COALESCE(
    (SELECT AVG(score)::DECIMAL(3,2) FROM ratings
     WHERE content_id = %s), 0)
WHERE content_id = %s;
\end{lstlisting}

\textbf{MySQL}
\begin{lstlisting}[language=SQL]
UPDATE content SET avg_rating = COALESCE(
    (SELECT AVG(score) FROM ratings
     WHERE content_id = %s), 0)
WHERE content_id = %s;
\end{lstlisting}

\textbf{MongoDB}
\begin{lstlisting}[language=Python]
db.ratings.aggregate([
    { $match: { content_id: <cid> } },
    { $group: { _id: null, avg: { $avg: "$score" } } }
]);

db.content.updateOne(
    { _id: <cid> },
    { $set: { avg_rating: <avg> } }
);
\end{lstlisting}

\textbf{Neo4j}
\begin{lstlisting}[language=SQL]
MATCH (:Profile)-[r:RATED]->(c:Content {content_id: $cid})
WITH c, avg(r.score) AS avgScore
SET c.avg_rating = avgScore;
\end{lstlisting}


\subsubsection{U3: Masowa zmiana planu subskrypcji}

Masowe przepisanie wszystkich aktywnych subskrypcji z planu Basic na Standard
w jednej operacji. Test bulk-update po filtrze bez klucza głównego -- pokazuje
zachowanie silnika przy równoczesnej modyfikacji wielu wierszy.

\textbf{PostgreSQL / MySQL}
\begin{lstlisting}[language=SQL]
UPDATE subscriptions
SET plan_name = 'Standard', price_monthly = 43.99,
    max_streams = 2, max_resolution = '1080p'
WHERE status = 'active' AND plan_name = 'Basic';
\end{lstlisting}

\textbf{MongoDB}
\begin{lstlisting}[language=Python]
db.users.updateMany(
    { "subscription.status": "active",
      "subscription.plan_name": "Basic" },
    { $set: {
        "subscription.plan_name": "Standard",
        "subscription.price_monthly": 43.99,
        "subscription.max_streams": 2,
        "subscription.max_resolution": "1080p"
    } }
);
\end{lstlisting}

\textbf{Neo4j}
\begin{lstlisting}[language=SQL]
MATCH (s:Subscription)
WHERE s.status = 'active' AND s.plan_name = 'Basic'
CALL (s) {
    SET s.plan_name = 'Standard', s.price_monthly = 43.99,
        s.max_streams = 2, s.max_resolution = '1080p'
} IN TRANSACTIONS OF 1000 ROWS;
\end{lstlisting}


\subsubsection{U4: Aktualizacja danych użytkownika}

Aktualizacja danych kontaktowych pojedynczego użytkownika (email + telefon +
updated\_at). Prosta aktualizacja jednego rekordu po user\_id -- przy wariancie z indeksami
unikalny indeks na email ma wpływ na koszt zapisu.

\textbf{PostgreSQL / MySQL}
\begin{lstlisting}[language=SQL]
UPDATE users
SET email = %s, phone = '+48 111 222 333', updated_at = NOW()
WHERE user_id = %s;
\end{lstlisting}

\textbf{MongoDB}
\begin{lstlisting}[language=Python]
db.users.updateOne(
    { _id: <uid> },
    { $set: {
        email: "updated_<uid>@bench.com",
        phone: "+48 111 222 333",
        updated_at: "2025-06-15 12:00:00"
    } }
);
\end{lstlisting}

\textbf{Neo4j}
\begin{lstlisting}[language=SQL]
MATCH (u:User {user_id: $uid})
SET u.email = $email;
\end{lstlisting}


\subsubsection{U5: Oznaczenie treści jako nieaktywnej}

Oznaczenie jednej treści jako nieaktywnej (is\_active = false) -- operacja
typu "soft delete". Minimalna aktualizacja jednego rekordu po kluczu głównym na kolumnie
często występującej w warunkach filtrów.

\textbf{PostgreSQL / MySQL}
\begin{lstlisting}[language=SQL]
UPDATE content SET is_active = FALSE WHERE content_id = %s;
\end{lstlisting}

\textbf{MongoDB}
\begin{lstlisting}[language=Python]
db.content.updateOne(
    { _id: <cid> },
    { $set: { is_active: false } }
);
\end{lstlisting}

\textbf{Neo4j}
\begin{lstlisting}[language=SQL]
MATCH (c:Content {content_id: $cid})
SET c.is_active = false;
\end{lstlisting}


\subsubsection{U6: Masowa aktualizacja wskaźnika popularności}

Przeliczenie popularity\_score dla wszystkich treści w bazie na podstawie
formuły: total\_views · 0.0001 + avg\_rating · 5 + liczba\_ocen · 0.1.
Typowa operacja batchowa -- testuje koszt agregacji (GROUP BY / aggregate /
pattern grafu) w połączeniu z masowym UPDATE po całej kolekcji.

\textbf{PostgreSQL}
\begin{lstlisting}[language=SQL]
UPDATE content
SET popularity_score = (
    total_views * 0.0001 +
    avg_rating * 5 +
    COALESCE(rc.cnt, 0) * 0.1
)
FROM (
    SELECT content_id, COUNT(*) AS cnt
    FROM ratings GROUP BY content_id
) rc
WHERE rc.content_id = content.content_id;
\end{lstlisting}

\textbf{MySQL}
\begin{lstlisting}[language=SQL]
UPDATE content c
JOIN (
    SELECT content_id, COUNT(*) AS cnt
    FROM ratings GROUP BY content_id
) rc ON rc.content_id = c.content_id
SET c.popularity_score = (
    c.total_views * 0.0001 +
    c.avg_rating * 5 +
    rc.cnt * 0.1
);
\end{lstlisting}

\textbf{MongoDB}
\begin{lstlisting}[language=Python]
db.ratings.aggregate([
    { $group: { _id: "$content_id", cnt: { $sum: 1 } } }
]);

db.content.bulkWrite([
    { updateOne: {
        filter: { _id: <cid> },
        update: { $set: { popularity_score:
            round(total_views*0.0001 + avg_rating*5 + cnt*0.1, 2) } }
    } },
    ...
], { ordered: false });
\end{lstlisting}

\textbf{Neo4j}
\begin{lstlisting}[language=SQL]
MATCH (c:Content)
CALL (c) {
    OPTIONAL MATCH (:Profile)-[r:RATED]->(c)
    WITH c, avg(r.score) AS avgScore, count(r) AS ratingCount
    SET c.popularity_score = (
        c.total_views * 0.0001 + c.avg_rating * 5 + ratingCount * 0.1
    )
} IN TRANSACTIONS OF 500 ROWS;
\end{lstlisting}


% -------------------------------------------------------
\subsection{DELETE (D1--D6)}
% -------------------------------------------------------

Grupa scenariuszy DELETE obejmuje sześć wariantów usuwania rekordów -- od lekkiego
usunięcia pojedynczego wpisu po kluczu złożonym, przez
usunięcie rekordu z pełną kaskadą zależności, aż po masowe czyszczenie danych po
 filtrze zakresu. Scenariusze pokazują różnice między mechanizmami kaskady
(\texttt{FK ON DELETE CASCADE} w SQL, aplikacyjne \texttt{deleteMany} w MongoDB,
\texttt{DETACH DELETE} w Neo4j) oraz wpływ indeksów na koszt wyszukiwania rekordów
przeznaczonych do usunięcia.

\subsubsection{D1: Usunięcie treści z kaskadą}

Usunięcie pojedynczej treści wraz z pełną kaskadą zależności: historia
oglądania, "moja lista", oceny, sezony, odcinki, obsada. W SQL dzieje się to
automatycznie dzięki FK ON DELETE CASCADE, w Mongo kaskadę wykonuje aplikacja
(4x deleteMany/deleteOne), w Neo4j użyto DETACH DELETE w dwóch krokach.

\textbf{PostgreSQL / MySQL}
\begin{lstlisting}[language=SQL]
DELETE FROM content WHERE content_id = %s;
-- kaskada na watch_history, my_list, ratings,
-- seasons, episodes, content_people przez FK ON DELETE CASCADE
\end{lstlisting}

\textbf{MongoDB}
\begin{lstlisting}[language=Python]
db.watch_history.deleteMany({ content_id: <cid> });
db.my_list.deleteMany({ content_id: <cid> });
db.ratings.deleteMany({ content_id: <cid> });
db.content.deleteOne({ _id: <cid> });
\end{lstlisting}

\textbf{Neo4j}
\begin{lstlisting}[language=SQL]
MATCH (c:Content {content_id: $cid})
OPTIONAL MATCH (c)-[:HAS_SEASON]->(s:Season)-[:HAS_EPISODE]->(e:Episode)
DETACH DELETE e, s;

MATCH (c:Content {content_id: $cid})
DETACH DELETE c;
\end{lstlisting}


\subsubsection{D2: Usunięcie profilu z historią}

Usunięcie pojedynczego profilu wraz z kaskadą wszystkich jego aktywności --
historia oglądania, "moja lista" oraz wystawione oceny. Scenariusz zbliżony do
D1, ale działa na bardziej rozproszonych relacjach (1 profil → N wpisów
w trzech różnych tabelach/kolekcjach).

\textbf{PostgreSQL / MySQL}
\begin{lstlisting}[language=SQL]
DELETE FROM profiles WHERE profile_id = %s;
-- kaskada: watch_history, my_list, ratings przez FK ON DELETE CASCADE
\end{lstlisting}

\textbf{MongoDB}
\begin{lstlisting}[language=Python]
db.watch_history.deleteMany({ profile_id: <pid> });
db.my_list.deleteMany({ profile_id: <pid> });
db.ratings.deleteMany({ profile_id: <pid> });
db.users.updateOne(
    { _id: <uid> },
    { $pull: { profiles: { profile_id: <pid> } } }
);
\end{lstlisting}

\textbf{Neo4j}
\begin{lstlisting}[language=SQL]
MATCH (p:Profile {profile_id: $pid})
DETACH DELETE p;
\end{lstlisting}


\subsubsection{D3: Czyszczenie starych ocen}

Masowe czyszczenie starych ocen -- usunięcie wszystkich rekordów ratings
sprzed daty odcięcia. Test bulk-delete po filtrze zakresu dat (created\_at <
cutoff). W Neo4j wymusza iteracyjne usuwanie po 10 000 krawędzi, żeby nie
przekroczyć pamięci transakcji.

\textbf{PostgreSQL / MySQL}
\begin{lstlisting}[language=SQL]
DELETE FROM ratings WHERE created_at < '2000-06-01';
\end{lstlisting}

\textbf{MongoDB}
\begin{lstlisting}[language=Python]
db.ratings.deleteMany({ created_at: { $lt: "2000-06-01" } });
\end{lstlisting}

\textbf{Neo4j}
\begin{lstlisting}[language=SQL]
-- petla po 10 000 rekordow do wyczerpania
MATCH ()-[r:RATED]->()
WHERE r.rated_at < $cutoff
WITH r LIMIT 10000
DELETE r
RETURN count(*) AS deleted;
\end{lstlisting}


\subsubsection{D4: Usunięcie pozycji z listy}

Usunięcie jednej pozycji z listy "Moja lista" wybranego profilu, po kluczu
złożonym (profile\_id, content\_id). Najlżejszy scenariusz DELETE w zestawie -- usunięcie jednego rekordu
po kluczu złożonym (profile\_id, content\_id).

\textbf{PostgreSQL / MySQL}
\begin{lstlisting}[language=SQL]
DELETE FROM my_list
WHERE profile_id = %s AND content_id = %s;
\end{lstlisting}

\textbf{MongoDB}
\begin{lstlisting}[language=Python]
db.my_list.deleteOne(
    { profile_id: <pid>, content_id: <cid> }
);
\end{lstlisting}

\textbf{Neo4j}
\begin{lstlisting}[language=SQL]
MATCH (p:Profile {profile_id: $pid})
      -[r:ADDED_TO_LIST]->
      (c:Content {content_id: $cid})
DELETE r;
\end{lstlisting}


\subsubsection{D5: Usunięcie subskrypcji z płatnościami}

Usunięcie pojedynczej subskrypcji wraz z kaskadą powiązanych płatności.
W SQL rozwiązane przez FK ON DELETE CASCADE, w Mongo przez aplikacyjne
deleteMany po subscription\_id, w Neo4j przez DETACH DELETE.

\textbf{PostgreSQL / MySQL}
\begin{lstlisting}[language=SQL]
DELETE FROM subscriptions WHERE subscription_id = %s;
-- kaskada na payments przez FK ON DELETE CASCADE
\end{lstlisting}

\textbf{MongoDB}
\begin{lstlisting}[language=Python]
db.payments.deleteMany({ subscription_id: <sid> });
\end{lstlisting}

\textbf{Neo4j}
\begin{lstlisting}[language=SQL]
MATCH (sub:Subscription {subscription_id: $sid})
OPTIONAL MATCH (sub)-[:HAS_PAYMENT]->(pay:Payment)
DETACH DELETE pay, sub;
\end{lstlisting}


\subsubsection{D6: Masowe usunięcie nieaktywnych użytkowników}

Masowe usunięcie użytkowników oznaczonych jako status = 'deleted' z podanego
zakresu ID (50 / 500 / 2000). Testuje bulk-delete po filtrze złożonym
(status + zakres user\_id) -- łączy odczyt po indeksie statusu z operacją
DELETE na wielu wierszach naraz.

\textbf{PostgreSQL / MySQL}
\begin{lstlisting}[language=SQL]
DELETE FROM users
WHERE status = 'deleted'
  AND user_id >= %s AND user_id < %s;
\end{lstlisting}

\textbf{MongoDB}
\begin{lstlisting}[language=Python]
db.users.deleteMany({
    status: "deleted",
    _id: { $gte: <start_uid>, $lt: <start_uid + n> }
});
\end{lstlisting}

\textbf{Neo4j}
\begin{lstlisting}[language=SQL]
MATCH (u:User)
WHERE u.status = 'deleted'
  AND u.user_id >= $start AND u.user_id < $end
DETACH DELETE u;
\end{lstlisting}